\documentclass[12pt]{report}
\usepackage[utf8]{inputenc}
\usepackage[T1]{fontenc}
\usepackage{amsmath}
\usepackage{amssymb}
\usepackage{amsthm}
\usepackage{geometry}
\usepackage{graphicx}
\usepackage{fancyhdr}
\usepackage{booktabs}
\setlength{\headheight}{15pt}
\usepackage[hidelinks]{hyperref}

\geometry{a4paper, margin=1in}
\parskip=0.5\baselineskip


\pagestyle{fancy}
\fancyhf{}
\fancyhead[L]{\rightmark}
\fancyhead[R]{\thepage}
\renewcommand{\headrulewidth}{0.4pt}
\fancypagestyle{plain}{
  \fancyhf{}
  \fancyhead[R]{\thepage}
  \renewcommand{\headrulewidth}{0pt}
}
\setcounter{secnumdepth}{3}

\begin{document}

\thispagestyle{empty}
\begin{center}
    \vspace*{1cm}

    \IfFileExists{sharif-university-logo.png}{
        \includegraphics[width=3cm]{sharif-university-logo.png}\\[0.5cm]
    }{}

    {\large \textbf{Sharif University of Technology}}
    
    {\large \textbf{Department of Physics}}\\[1.5cm]

    \vspace{1.5cm}

    {\Huge \textbf{QFT 1: Problem Set 1 }}\\[1.5cm]


    \begin{tabular}{p{3cm} p{8cm}}
        \textbf{Professor:} & Arash Arabi Ardehali  \\[0.5cm]
        \textbf{Student:} & Abolfazl Balouch \\[0.5cm]
        \textbf{E-mail:} & Balouch1381@gmail.com \\[0.5cm]
    \end{tabular}
    \vfill
    {\large \today}
\end{center}
\newpage

\section*{Problem}
Consider a 4-dimensional hyper-cubic lattice with the lattice spacing \(\Delta\). Compute the ``lattice propagator'' as the sum over all lattice paths connecting the lattice points 0 and \(x\),
\[
\sum_{\text{lattice path } 0 \to x} e^{-m_0 L(\text{path})},
\]
where \(L(\text{path})\) is the number of links in the path. The path may go through a given link any number of times.

Consider the limit \(\Delta \to 0\). Observe that at fixed \(m_0\) the sum either diverges, or is dominated by the shortest path from 0 to \(x\). Determine how \(m_0 = m_0(\Delta)\) should depend on \(\Delta\) for long paths to remain relevant in the limit \(\Delta \to 0\).

Suggestion: follow the steps outlined in Lecture 7, with suitable modifications to take into account the lattice structure. The representation
\[
\delta_{n,0} = \int_{-\pi}^{\pi} \frac{dk}{2\pi} e^{ikn}
\]
of the Kronecker delta may be useful.

\section*{Solution}
Suppose point $x$ is on the lattice and $x=\Delta N$ ($N=(N_1,...,N_4)$ is the integer coordinates of the lattice. meaning that $N_\mu=\frac{x_\mu}{\Delta}$.) If \(C_\Delta(N)\) is the number of \(n\)-step paths from \(0\) to \(N\), we have:

\begin{equation}
G_\Delta(N) = \sum_{n=0}^{\infty} e^{-m_0 n} C_n(N).
\label{eq1}
\end{equation}

In four dimensions, $\delta_{n,0} = \int_{-\pi}^{\pi} \frac{dk}{2\pi} e^{ikn}$ takes the form $\delta_{N,0} = \int_{-\pi}^{\pi} \frac{d^4k}{(2\pi)^4} e^{ik \cdot N}$. In one step, we can move in both the \(\hat{e}_\mu\)-direction and the \(-\hat{e}_\mu\)-direction, so the coefficient corresponding to one step in Fourier space is equal to \(e^{ik_\mu} + e^{-ik_\mu} \). The sum over the four directions is equal to

\begin{equation}
\sum_{\mu=1}^{4} \left( e^{i k \mu} + e^{-i k \mu} \right) = 2 \sum_{\mu=1}^{4} \cos k_\mu
\label{eq2}
\end{equation}

Therefore, for \(n\) steps we have:

\begin{equation}
C_n(N) = \int_{-\pi}^{\pi} \frac{d^4 k}{(2\pi)^4}\,
e^{-i k \cdot N}
\left( 2 \sum_{\mu=1}^{4} \cos k_\mu \right)^n
\label{eq3}
\end{equation}

Substituting Equation \eqref{eq3} into Equation \eqref{eq1}, we obtain:

\begin{equation}
G_\Delta(N)
=
\int_{-\pi}^{\pi}
\frac{d^4k}{(2\pi)^4}\,
e^{-i k\cdot N}
\sum_{n=0}^{\infty}
\left[
e^{-m_0}\,
2 \sum_{\mu=1}^{4} \cos k_\mu
\right]^n
\label{eq4}
\end{equation}

We use the geometric series.($\sum_{n=0}^{\infty} r^n = \frac{1}{1-r}$)

\begin{equation}
G_\Delta(N)
=
\int_{-\pi}^{\pi}
\frac{d^4 k}{(2\pi)^4}\,
\frac{e^{-i k\cdot N}}
{1 - 2 e^{-m_0} \sum_{\mu=1}^{4} \cos k_\mu}
\label{eq5}
\end{equation}

To see continuum physics, we need to look at small momenta (\(k_\mu \ll 1\) because $p_\mu=\frac{k_\mu}{\Delta}$ and $p_\mu$ is finite therefor when $ \Delta \rightarrow 0 \Rightarrow k_\mu  \rightarrow 0$):

\begin{equation}
\cos k_\mu = 1 - \frac{k_\mu^2}{2} + \frac{k_\mu^4}{24} + \cdots \Rightarrow 2\cos k_\mu = 2 - k_\mu^2 + O(k_\mu^4)
\label{eq6}
\end{equation}

Summing over the four dimensions:

\begin{equation}
2\sum_{\mu=1}^4 \cos k_\mu = 8 - k^2 + O(k^4) \quad ,(k^2 = \sum_{\mu=1}^4 k_\mu^2)
\label{eq7}
\end{equation}

So the denominator of Equation \eqref{eq5} becomes as follows:

\begin{equation}
1 - 8e^{-m_0} + e^{-m_0} k^2 + O(k^4)
\label{eq8}
\end{equation}

We need to bring the denominator into the familiar form $\frac{1}{p^2+m^2}$. ($p_\mu=\frac{k_\mu}{\Delta} \Rightarrow k_\mu=\Delta p_\mu \Rightarrow k^2=\Delta^2p^2$) Also, from Equation \eqref{eq8} we see that the coefficient of \(p^2\) must be \(e^{-m_0}\Delta\); therefore:

\begin{equation}
\begin{split}
1-8e^{-m_0}+e^{-m_0}\Delta^2p^2
=
e^{-m_0}\Delta^2(p^2+m^2) \Rightarrow 1-8e^{-m_0}=e^{-m_0}m^2\Delta^2 \\
\Rightarrow e^{m_0}-8=m^2\Delta^2 \Rightarrow e^{m_0}=8+m^2\Delta^2 \Rightarrow \boxed{
m_0(\Delta)=\ln(8+m^2\Delta^2).
}
\label{eq9}
\end{split} 
\end{equation}

We write Equation \eqref{eq9} in the following form:

\begin{equation}
m_0 = \ln 8 + \ln\left(1+\frac{m^2\Delta^2}{8}\right)
\label{eq10}
\end{equation}

For $\Delta \rightarrow 0$ Using $ \ln(1+\epsilon)=\epsilon-\frac{\epsilon^2}{2}+\cdots$, we thus obtain:

\begin{equation}
\boxed{\boxed{
m_0(\Delta)
=
\ln 8+\frac{m^2}{8}\Delta^2+O(\Delta^4)
}}
\label{eq11}
\end{equation}

\newpage
\subsection*{AI use declaration}
Model(s) used: 

ChatGPT (GPT-5.6)\\

Revisions/corrections made to model output, and clarifications added:

I didn't give the full question to the AI, but I used it in the solution process. For writing the LaTeX code, the parts I used from it in the conceptual process of the question are $K_\mu=p_\mu/\Delta$ and generally concepts related to lattices.\\

Comparisons/checks against the lecture notes:

Before reading the question, I thoroughly reviewed the lecture notes, performed all intermediate calculations, and wrote the question based on that.\\

Calculation steps, limiting cases, or consistency checks verified by hand:

Since this problem did not involve complex calculations, I handled all the calculations myself.

Calculation steps, limiting cases, or consistency checks subjected to automated verification (via Mathematica, python, LLMs):

As I said above I didn't give the full question to the AI, but I used it in the solution process. For writing the LaTeX code, the parts I used from it in the conceptual process of the question are $K_\mu=p_\mu/\Delta$ and generally concepts related to lattices.\\





\end{document}
